\chapter{Review of the Great Theorems}

\section{Hahn--Banach theory}

\begin{notetheorem}[Hahn--Banach theorem over real vector spaces]
Let $X$ be a real vector space, let $Y\subseteq X$ be a linear subspace, and
let $p\colon X\to\mathbb R$ be sublinear:
\[
  p(\alpha x)=\alpha p(x)\quad(\alpha\ge0),
  \qquad
  p(x+y)\le p(x)+p(y).
\]
If $f\colon Y\to\mathbb R$ is linear and $f(y)\le p(y)$ for every $y\in Y$,
then there is a linear functional $\widetilde f\colon X\to\mathbb R$ such that
\[
  \widetilde f|_Y=f,
  \qquad
  \widetilde f(x)\le p(x)\quad(x\in X).
\]
\end{notetheorem}

\begin{proof}
We follow the maximal-extension argument indicated by the Zorn-lemma
discussion in the source.  Let $\mathscr P$ be the set of pairs $(Z,g)$ such
that
\[
  Y\subseteq Z\subseteq X,\qquad
  Z\text{ is a linear subspace},\qquad
  g\colon Z\to\mathbb R\text{ is linear},
\]
and
\[
  g|_Y=f,\qquad g(z)\le p(z)\quad(z\in Z).
\]
Order $\mathscr P$ by extension:
\[
  (Z_1,g_1)\preccurlyeq(Z_2,g_2)
  \quad\Longleftrightarrow\quad
  Z_1\subseteq Z_2\ \text{ and }\ g_2|_{Z_1}=g_1.
\]
The pair $(Y,f)$ belongs to $\mathscr P$, so $\mathscr P$ is nonempty.

Let $\mathscr C\subseteq\mathscr P$ be a chain and set
\[
  Z_{\mathscr C}=\bigcup_{(Z,g)\in\mathscr C}Z.
\]
Because the members of the chain are totally ordered by inclusion,
$Z_{\mathscr C}$ is a linear subspace.  For $z\in Z_{\mathscr C}$ choose
$(Z,g)\in\mathscr C$ with $z\in Z$ and define
\[
  g_{\mathscr C}(z)=g(z).
\]
This is well defined: if $z$ lies in two members of the chain, one extends
the other and the two functionals agree on the smaller domain.  The same
chain argument shows that $g_{\mathscr C}$ is linear, extends $f$, and
satisfies $g_{\mathscr C}\le p$.  Thus every chain has an upper bound, and
Zorn's lemma gives a maximal pair $(M,g)$.

Suppose that $M\ne X$ and choose $x_0\in X\setminus M$.  Every linear
extension of $g$ to $M+\mathbb R x_0$ has the form
\[
  g_\alpha(m+t x_0)=g(m)+t\alpha
  \qquad(m\in M,\ t\in\mathbb R)
\]
for some $\alpha\in\mathbb R$.  Define
\[
  L=\sup_{m\in M}\bigl(g(m)-p(m-x_0)\bigr),
  \qquad
  U=\inf_{m\in M}\bigl(p(m+x_0)-g(m)\bigr).
\]
For arbitrary $m,n\in M$,
\[
\begin{aligned}
  g(m)+g(n)
   &=g(m+n)
    \le p(m+n)\\
   &\le p(m-x_0)+p(n+x_0),
\end{aligned}
\]
so
\[
  g(m)-p(m-x_0)\le p(n+x_0)-g(n).
\]
Taking the supremum in $m$ and the infimum in $n$ gives $L\le U$.  Choose
$\alpha\in[L,U]$.

If $t>0$, then, with $m/t\in M$,
\[
\begin{aligned}
  g_\alpha(m+t x_0)
   &=t\left(g(m/t)+\alpha\right)\\
   &\le t\left(g(m/t)+p(m/t+x_0)-g(m/t)\right)\\
   &=p(m+t x_0).
\end{aligned}
\]
If $t<0$, write $t=-s$ with $s>0$.  Since $\alpha\ge L$,
\[
\begin{aligned}
  g_\alpha(m-sx_0)
   &=s\left(g(m/s)-\alpha\right)\\
   &\le s\,p(m/s-x_0)
    =p(m-sx_0).
\end{aligned}
\]
The case $t=0$ is the original domination on $M$.  Thus
$(M+\mathbb R x_0,g_\alpha)$ is a proper extension of $(M,g)$, contradicting
maximality.  Therefore $M=X$ and $g$ is the required extension.
\end{proof}

\begin{notecorollary}[Norm-preserving Hahn--Banach extension]
Let $X$ be a normed space, let $Y\subseteq X$ be a linear subspace, and let
$f\in Y^*$.  There is $\widetilde f\in X^*$ satisfying
\[
  \widetilde f|_Y=f,
  \qquad
  \|\widetilde f\|=\|f\|.
\]
\end{notecorollary}

\begin{proof}
If the scalar field is real, apply the preceding theorem with
$p(x)=\|f\|\,\|x\|$.  Since this sublinear function is symmetric, domination
of both $x$ and $-x$ gives
\[
  |\widetilde f(x)|\le \|f\|\,\|x\|,
\]
and therefore $\|\widetilde f\|\le\|f\|$.  Restriction to $Y$ gives the
reverse inequality.

For completeness, suppose that $X$ is complex.  Regard $X$ and $Y$ as real
spaces and put $h=\operatorname{Re}f$.  The real theorem gives a real-linear
extension $\widetilde h$ satisfying
\[
  |\widetilde h(x)|\le\|f\|\,\|x\|.
\]
Define
\[
  \widetilde f(x)=\widetilde h(x)-i\widetilde h(ix).
\]
Then $\widetilde f(ix)=i\widetilde f(x)$, so $\widetilde f$ is complex linear,
and on $Y$ one has
\[
  \widetilde f(y)
  =\operatorname{Re}f(y)-i\operatorname{Re}f(iy)
  =f(y).
\]
For $|\zeta|=1$,
\[
  \operatorname{Re}\bigl(\zeta\widetilde f(x)\bigr)
  =\widetilde h(\zeta x)
  \le\|f\|\,\|x\|.
\]
Taking the supremum over unimodular $\zeta$ yields
$|\widetilde f(x)|\le\|f\|\,\|x\|$.  Again restriction gives equality of
the norms.
\end{proof}

\begin{noteremark}[A useful uniqueness criterion]
If $X^*$ is strictly convex, then every norm-preserving Hahn--Banach extension
is unique.  Indeed, if $F_1$ and $F_2$ are two distinct extensions of $f$ with
$\|F_1\|=\|F_2\|=\|f\|>0$, then $(F_1+F_2)/2$ is another extension and must
have norm at least $\|f\|$; this contradicts strict convexity after
normalization.  The converse is not automatic from the one-line argument in
the source and is therefore not asserted here without an additional global
unique-extension hypothesis.
\end{noteremark}

\begin{noteproposition}[Norming functionals]
Let $X$ be a normed space and let $x\in X\setminus\{0\}$.  There is
$f_x\in X^*$ such that
\[
  f_x(x)=\|x\|,
  \qquad
  \|f_x\|=1.
\]
Consequently, continuous linear functionals separate points: if $x\ne y$,
then some $f\in X^*$ satisfies $f(x)\ne f(y)$.
\end{noteproposition}

\begin{proof}
Define $g(tx)=t\|x\|$ on $\operatorname{span}\{x\}$.  Then $\|g\|=1$ and a
norm-preserving Hahn--Banach extension gives $f_x$.  Apply this to $x-y$ for
the separation statement.
\end{proof}

\section{Geometric separation}

The order inequalities in this section are stated for real normed spaces.  In
a complex normed space one applies the result to the underlying real space and
then writes the resulting bounded real-linear functional $h$ as
\[
  h=\operatorname{Re}f,
  \qquad
  f(x)=h(x)-ih(ix),
\]
with $f$ complex linear.  Thus the corresponding complex statements use
$\operatorname{Re}f$ in place of $f$.

\begin{notelemma}[Minkowski functional]
Let $X$ be a normed space and let $C\subseteq X$ be open and convex with
$0\in C$.  Define
\[
  p_C(x)=\inf\{t>0:x\in tC\}.
\]
Then $p_C$ is finite and sublinear.  Moreover, for some $M>0$,
\[
  p_C(x)\le M\|x\|\quad(x\in X),
  \qquad
  C=\{x\in X:p_C(x)<1\}.
\]
\end{notelemma}

\begin{proof}
Because $C$ is open and contains $0$, there is $r>0$ such that
$B(0,r)\subseteq C$.  For $x\ne0$ and every $t>\|x\|/r$, we have
$\|x/t\|<r$, hence $x/t\in C$ and therefore $x\in tC$.  Consequently,
\[
  0\le p_C(x)\le \frac{\|x\|}{r},
\]
and the same inequality is trivial for $x=0$.

For $\lambda>0$,
\[
\begin{aligned}
  p_C(\lambda x)
   &=\inf\{t>0:\lambda x\in tC\}\\
   &=\inf\{\lambda s:s>0,\ x\in sC\}
    =\lambda p_C(x).
\end{aligned}
\]
Also $p_C(0)=0$, so positive homogeneity holds for every $\lambda\ge0$.

To prove subadditivity, fix $\varepsilon>0$.  By the definition of the
infimum choose $a,b>0$ such that
\[
  x\in aC,\quad y\in bC,\quad
  a<p_C(x)+\varepsilon,\quad b<p_C(y)+\varepsilon.
\]
Then $x/a,y/b\in C$, and convexity gives
\[
  \frac{x+y}{a+b}
   =\frac{a}{a+b}\frac{x}{a}
    +\frac{b}{a+b}\frac{y}{b}\in C.
\]
Thus
\[
  p_C(x+y)\le a+b
  <p_C(x)+p_C(y)+2\varepsilon.
\]
Letting $\varepsilon\downarrow0$ proves subadditivity.

Finally, if $x\in C$, continuity of $s\mapsto sx$ and openness of $C$ give
$s>1$ sufficiently close to $1$ with $sx\in C$.  Hence
$x\in s^{-1}C$ and $p_C(x)\le s^{-1}<1$.  Conversely, if $p_C(x)<1$, choose
$t\in(p_C(x),1)$ with $x\in tC$.  Writing $x=tc$ with $c\in C$, convexity
of $C$ and $0\in C$ imply $tc\in C$.  Therefore
\[
  C=\{x\in X:p_C(x)<1\}.
\]
\end{proof}

\begin{notelemma}[Separation of an open convex set and a point]
Let $X$ be a real normed space, let $C\subseteq X$ be open and convex, and let
$x_0\notin C$.  Then there are
$f\in X^*\setminus\{0\}$ and $\alpha\in\mathbb R$ such that
\[
  f(x)<\alpha=f(x_0)\qquad(x\in C).
\]
\end{notelemma}

\begin{proof}
Choose $c_0\in C$ and replace $C$ by $C-c_0$ and $x_0$ by $x_0-c_0$.
After this translation, $0\in C$ and still $x_0\notin C$.  Let $p_C$ be the
Minkowski functional of the translated set.  Since $C=\{p_C<1\}$,
$p_C(x_0)\ge1$.

Define
\[
  g(t x_0)=t
  \qquad(t\in\mathbb R)
\]
on $\operatorname{span}\{x_0\}$.  For $t\ge0$,
\[
  p_C(t x_0)=t p_C(x_0)\ge t=g(t x_0),
\]
whereas for $t<0$ we have $g(t x_0)<0\le p_C(t x_0)$.  Thus $g\le p_C$ on
its domain.  Hahn--Banach gives a linear extension $f$ with $f\le p_C$ on
$X$ and $f(x_0)=1$.

The norm estimate for $p_C$ also proves continuity: for every $x\in X$,
\[
  |f(x)|\le\max\{p_C(x),p_C(-x)\}\le M'\|x\|.
\]
Hence $f\in X^*$.  If $x\in C$, then
\[
  f(x)\le p_C(x)<1=f(x_0).
\]
Undoing the translation gives the required strict separation, with
$\alpha=f(x_0)$.
\end{proof}

\begin{notetheorem}[Geometric Hahn--Banach separation]
Let $A,B$ be nonempty disjoint convex subsets of a real normed space $X$.

\begin{enumerate}
  \item If $A$ is open, then there are $f\in X^*\setminus\{0\}$ and
        $\alpha\in\mathbb R$ such that
        \[
          f(a)<\alpha\le f(b)
          \qquad(a\in A,\ b\in B).
        \]
  \item If $A$ is closed and $B$ is compact, then there are
        $f\in X^*\setminus\{0\}$, $\alpha\in\mathbb R$, and $\varepsilon>0$
        such that
        \[
          f(a)+\varepsilon\le\alpha\le f(b)
          \qquad(a\in A,\ b\in B).
        \]
\end{enumerate}
\end{notetheorem}

\begin{proof}
For the first assertion, put
\[
  C=A-B=\{a-b:a\in A,\ b\in B\}.
\]
The set $C$ is convex and open, and $0\notin C$.  Apply the preceding lemma
to $C$ and the point $0$.  After changing the sign of the functional if
necessary, we obtain $f\in X^*\setminus\{0\}$ such that
\[
  f(a-b)<0
  \qquad(a\in A,\ b\in B).
\]
Hence $f(a)<f(b)$ for all $a\in A$ and $b\in B$, and
\[
  -\infty<\sup_{a\in A}f(a)
  \le\inf_{b\in B}f(b)<\infty.
\]
Let $\alpha=\sup_A f$.  Since $A$ is open, no point of $A$ can maximize the
nonzero functional $f$: for $a\in A$, choose $h$ with $f(h)>0$; then
$a+th\in A$ for sufficiently small $t>0$ and $f(a+th)>f(a)$.  Therefore
\[
  f(a)<\alpha\le f(b)
  \qquad(a\in A,\ b\in B).
\]

For the second assertion, $A-B$ is closed.  Indeed, if
$a_n-b_n\to z$, compactness of $B$ gives a subsequence $b_{n_k}\to b\in B$;
then $a_{n_k}\to z+b\in A$, because $A$ is closed, and so $z\in A-B$.
Since $0\notin A-B$, choose $r>0$ with
\[
  B(0,r)\cap(A-B)=\varnothing.
\]
Apply the first part to the open convex set $B(0,r)$ and the convex set
$A-B$.  There are $g\in X^*\setminus\{0\}$ and $\beta\in\mathbb R$ such
that
\[
  g(u)<\beta\le g(a-b)
  \qquad(\|u\|<r,\ a\in A,\ b\in B).
\]
Taking the supremum over the open ball gives
\[
  r\|g\|\le g(a-b)
  \qquad(a\in A,\ b\in B).
\]
Set $f=-g$ and $\delta=r\|g\|>0$.  Then
\[
  f(a)+\delta\le f(b)
  \qquad(a\in A,\ b\in B),
\]
so $\sup_Af+\delta\le\inf_Bf$.  With
\[
  \varepsilon=\frac\delta2,
  \qquad
  \alpha=\frac12\left(\sup_Af+\inf_Bf\right),
\]
we obtain
\[
  f(a)+\varepsilon\le\alpha\le f(b)
  \qquad(a\in A,\ b\in B).
\]
\end{proof}

\section{Baire category and uniform boundedness}

\begin{notedefinition}
A subset $N$ of a topological space is \emph{nowhere dense} if
$\operatorname{int}(\overline N)=\varnothing$.  A space is of the
\emph{first category} if it is a countable union of nowhere dense sets, and of
the \emph{second category} otherwise.
\end{notedefinition}

\begin{notelemma}
For a subset $A$ of a topological space $X$, the following are equivalent:
\begin{enumerate}
  \item $A$ is dense in $X$;
  \item every nonempty open subset of $X$ meets $A$;
  \item $\operatorname{int}(X\setminus A)=\varnothing$.
\end{enumerate}
\end{notelemma}

\begin{proof}
Assume first that $\overline A=X$, and let $O\subseteq X$ be nonempty and
open.  Choose $x\in O$.  Since $x\in\overline A$, every neighbourhood of
$x$, in particular $O$, meets $A$.  Thus (1) implies (2).

If (2) holds and $\operatorname{int}(X\setminus A)$ were nonempty, that
interior would be a nonempty open set disjoint from $A$, contradicting (2).
Hence (2) implies (3).

Finally, the elementary identity
\[
  X\setminus\overline A=\operatorname{int}(X\setminus A)
\]
shows that (3) is equivalent to $X\setminus\overline A=\varnothing$, that is,
$\overline A=X$.  This proves (3) implies (1) and completes the equivalence.
\end{proof}

\begin{notetheorem}[Baire category theorem]
Every nonempty complete metric space is of the second category.  Equivalently:
\begin{enumerate}
  \item a countable union of closed sets with empty interior has empty
        interior;
  \item a countable intersection of open dense sets is dense.
\end{enumerate}
\end{notetheorem}

\begin{proof}
We prove the open-dense formulation, following the nested-ball construction
in the source but using closed balls so that the limiting point is retained.
Let $(U_n)_{n\ge1}$ be open and dense, and let $O\subseteq X$ be nonempty
and open.  Since $O\cap U_1$ is nonempty and open, choose $x_1\in X$ and
$0<r_1<1$ such that
\[
  \overline B(x_1,r_1)\subseteq O\cap U_1.
\]
Inductively, once $\overline B(x_n,r_n)$ has been chosen, the set
$B(x_n,r_n)\cap U_{n+1}$ is nonempty and open.  Choose
$x_{n+1}\in X$ and
\[
  0<r_{n+1}<\min\left\{\frac{r_n}{2},\frac1{n+1}\right\}
\]
with
\[
  \overline B(x_{n+1},r_{n+1})
  \subseteq B(x_n,r_n)\cap U_{n+1}.
\]
The closed balls are nested and $r_n\to0$.  If $m>n$, then
$x_m\in\overline B(x_n,r_n)$, so $d(x_m,x_n)\le r_n$.  Thus $(x_n)$ is
Cauchy.  Completeness gives $x_n\to x\in X$.

For every fixed $n$, all $x_m$ with $m\ge n$ lie in the closed ball
$\overline B(x_n,r_n)$, hence
\[
  x\in\overline B(x_n,r_n)\subseteq U_n.
\]
Also $x\in\overline B(x_1,r_1)\subseteq O$.  Consequently,
\[
  O\cap\bigcap_{n=1}^{\infty}U_n\ne\varnothing.
\]
Since every nonempty open set meets the intersection, $\bigcap_nU_n$ is
dense.

If $F_n$ is closed with empty interior, then $U_n=X\setminus F_n$ is open
and dense.  The preceding result says that $\bigcap_nU_n$ is dense, so its
complement $\bigcup_nF_n$ has empty interior.  This is the closed-set
formulation and proves the category theorem.
\end{proof}

\begin{noteexample}
A normed space with a countably infinite Hamel basis cannot be complete: it is
the union of the finite-dimensional spans of the first $n$ basis vectors, each
of which is closed and has empty interior.  Baire's theorem also underlies the
existence of continuous nowhere differentiable functions on $[0,1]$.
\end{noteexample}

\begin{notetheorem}[Uniform boundedness principle]
Let $X$ be a Banach space, let $Y$ be a normed space, and let
$\{T_i:i\in I\}\subseteq\Bop{X}{Y}$.  If
\[
  \sup_{i\in I}\|T_i x\|<\infty
  \qquad(x\in X),
\]
then
\[
  \sup_{i\in I}\|T_i\|<\infty.
\]
\end{notetheorem}

\begin{proof}
For $n\in\mathbb N$, define
\[
  F_n
  =\{x\in X:\|T_i x\|\le n\text{ for every }i\in I\}
  =\bigcap_{i\in I}T_i^{-1}\bigl(\overline B_Y(0,n)\bigr).
\]
Each $F_n$ is closed.  Pointwise boundedness says that every $x\in X$ lies
in some $F_n$, so $X=\bigcup_nF_n$.  Baire's theorem gives
$N\in\mathbb N$, $x_0\in X$, and $r>0$ such that
\[
  B(x_0,r)\subseteq F_N.
\]

To move the estimate to the origin, take $h\in X$ with $\|h\|<r$.  Both
$x_0+h$ and $x_0-h$ lie in $F_N$, and therefore, for every $i\in I$,
\[
\begin{aligned}
  2\|T_i h\|
   &=\|T_i(x_0+h)-T_i(x_0-h)\|\\
   &\le\|T_i(x_0+h)\|+\|T_i(x_0-h)\|
    \le2N.
\end{aligned}
\]
Hence $\|T_i h\|\le N$ for all $\|h\|<r$, uniformly in $i$.
If $\|x\|\le1$, put $h=(r/2)x$.  Then
\[
  \|T_i x\|=\frac2r\|T_i h\|\le\frac{2N}{r}.
\]
Taking the supremum over the unit ball gives
\[
  \sup_{i\in I}\|T_i\|\le\frac{2N}{r}<\infty.
\]
\end{proof}

\begin{notetheorem}[Banach--Steinhaus pointwise-limit theorem]
Let $X$ be Banach, let $Y$ be normed, and let
$T_n\in\Bop{X}{Y}$.  Suppose $T_nx$ converges for every $x\in X$, and define
$Tx=\lim_nT_nx$.  Then $T\in\Bop{X}{Y}$ and
\[
  \|T\|\le\liminf_{n\to\infty}\|T_n\|.
\]
\end{notetheorem}

\begin{proof}
Linearity is inherited from the pointwise limit.  For every $x$, the sequence
$(T_nx)$ is bounded, so uniform boundedness yields
$M:=\sup_n\|T_n\|<\infty$.  Hence $\|Tx\|\le M\|x\|$.  More precisely,
for $\|x\|=1$,
\[
  \|Tx\|=\lim_n\|T_nx\|
  \le\liminf_n\|T_n\|,
\]
and taking the supremum over the unit sphere gives the claim.
\end{proof}

\begin{notetheorem}[Weak boundedness implies norm boundedness]
Let $B\subseteq X$, where $X$ is a normed space.  If $f(B)$ is bounded in
$\mathbb K$ for every $f\in X^*$, then $B$ is norm bounded.
\end{notetheorem}

\begin{proof}
For $b\in B$, define $J_b\in(X^*)^*$ by $J_b(f)=f(b)$.  The family
$\{J_b:b\in B\}$ is pointwise bounded on the Banach space $X^*$, so uniform
boundedness gives $\sup_{b\in B}\|J_b\|<\infty$.  Hahn--Banach gives
$\|J_b\|=\|b\|$, and the result follows.
\end{proof}

\section{Open mapping, bounded inverse, and equivalent norms}

\begin{notedefinition}
A map between topological spaces is \emph{open} if it sends open sets to open
sets.
\end{notedefinition}

\begin{notetheorem}[Banach open mapping theorem]
Let $X$ and $Y$ be Banach spaces and let $T\in\Bop{X}{Y}$ be surjective.  Then
$T$ is an open map.
\end{notetheorem}

\begin{proof}
We retain the four steps of the handwritten proof.

\emph{Step 1: obtain interior in the closure of an image.}
Surjectivity gives
\[
  Y=\bigcup_{n=1}^{\infty}T(nB_X)
   \subseteq\bigcup_{n=1}^{\infty}\overline{T(nB_X)}.
\]
Thus the closed sets $\overline{T(nB_X)}$ cover the Banach space $Y$.
Baire's theorem yields $n_0\in\mathbb N$, $y_0\in Y$, and $r>0$ such that
\[
  B_Y(y_0,r)\subseteq\overline{T(n_0B_X)}.
\]

\emph{Step 2: center the ball at the origin.}
Let $\|z\|<r$.  Since $y_0$ and $y_0+z$ both lie in the displayed ball,
choose $x_k,u_k\in n_0B_X$ with
\[
  Tx_k\to y_0+z,
  \qquad
  Tu_k\to y_0.
\]
Then $T(x_k-u_k)\to z$ and $x_k-u_k\in2n_0B_X$.  Hence
\[
  B_Y(0,r)\subseteq\overline{T(2n_0B_X)}.
\]
After scaling, there is $\delta>0$ such that
\[
  B_Y(0,\delta)\subseteq\overline{T(B_X)}.
\]
For the correction argument we normalize this to
\[
  B_Y(0,1)\subseteq\overline{T(B_X)}.
\]

\emph{Step 3: remove the closure by successive corrections.}
Let $\|y\|<1/2$ and set $r_0=y$.  We construct $x_k\in X$ such that
\[
  \|x_k\|<2^{-k},
  \qquad
  r_k=y-T\sum_{j=1}^kx_j,
  \qquad
  \|r_k\|<2^{-k-1}.
\]
Suppose $r_{k-1}$ is known and $\|r_{k-1}\|<2^{-k}$.  Then
$2^kr_{k-1}\in B_Y(0,1)$.  By the closure inclusion, choose $u_k\in B_X$
such that
\[
  \|2^kr_{k-1}-Tu_k\|<\frac12.
\]
Set $x_k=2^{-k}u_k$.  Then $\|x_k\|<2^{-k}$ and
\[
  \|r_{k-1}-Tx_k\|<2^{-k-1},
\]
which is exactly the desired estimate for $r_k$.

Because $\sum_k\|x_k\|<1$, the series $\sum_kx_k$ converges in $X$ to some
$x\in B_X$.  Moreover, $r_k\to0$, so continuity of $T$ gives
\[
  Tx=T\left(\sum_{k=1}^{\infty}x_k\right)
     =\lim_{k\to\infty}T\sum_{j=1}^kx_j
     =y.
\]
Therefore $B_Y(0,1/2)\subseteq T(B_X)$.  Undoing the normalization, there is
$c>0$ such that
\[
  B_Y(0,c)\subseteq T(B_X).
\]

\emph{Step 4: pass to an arbitrary open set.}
Let $O\subseteq X$ be open and let $y=Tx\in T(O)$.  Choose
$\varepsilon>0$ with $x+\varepsilon B_X\subseteq O$.  Then
\[
  T(x+\varepsilon B_X)
  =y+\varepsilon T(B_X)
  \supseteq B_Y(y,\varepsilon c).
\]
Thus every point of $T(O)$ is an interior point, and $T(O)$ is open.
\end{proof}

\begin{notecorollary}[Bounded inverse theorem]
A bounded bijective linear map between Banach spaces has a bounded inverse.
\end{notecorollary}

\begin{proof}
The open mapping theorem says that $T^{-1}$ is continuous at $0$, hence
bounded.
\end{proof}

\begin{noteexample}[A priori estimate for a boundary-value problem]
Let
\[
 X=\{u\in C^2[0,1]:u(0)=u(1)=0\},
 \quad
 \|u\|_X=\|u\|_\infty+\|u'\|_\infty+\|u''\|_\infty,
\]
and let $Y=C[0,1]$ with the supremum norm.  For continuous coefficients
$a,b,c$, the operator
\[
  Lu=a u''+b u'+cu
\]
is bounded from $X$ to $Y$.  If the boundary-value problem $Lu=f$ has a
unique solution $u\in X$ for every $f\in Y$---equivalently, if $L$ is
bijective---then the bounded inverse theorem gives a constant $C$ such that
\[
  \|u\|_X\le C\|Lu\|_\infty.
\]
The bijectivity hypothesis is essential; it does not follow merely from the
formula for $L$.
\end{noteexample}

\begin{notetheorem}[Equivalence of complete norms]
Let $\|\cdot\|_1$ and $\|\cdot\|_2$ be norms making the same vector space $X$
Banach.  If
\[
  \|x\|_2\le C\|x\|_1\qquad(x\in X),
\]
then there is $C'>0$ such that
\[
  \|x\|_1\le C'\|x\|_2\qquad(x\in X).
\]
Thus the two norms are equivalent and induce the same topology.
\end{notetheorem}

\begin{proof}
The identity map $(X,\|\cdot\|_1)\to(X,\|\cdot\|_2)$ is a bounded bijection
between Banach spaces.  Apply the bounded inverse theorem.
\end{proof}

\section{Closed operators and the closed graph theorem}

\begin{notedefinition}
For a linear operator $T\colon\Dom(T)\subseteq X\to Y$, its graph is
\[
  \Graph(T)=\{(x,Tx):x\in\Dom(T)\}\subseteq X\times Y.
\]
The operator is \emph{closed} if its graph is closed.  Equivalently,
whenever $x_n\in\Dom(T)$, $x_n\to x$ in $X$, and $Tx_n\to y$ in $Y$, one
has $x\in\Dom(T)$ and $Tx=y$.
\end{notedefinition}

\begin{noteremark}
Every bounded linear operator is closed.  The converse fails when the domain
is not complete.  For example,
\[
  D\colon C^1[0,1]\subset C[0,1]\to C[0,1],
  \qquad Du=u',
\]
is closed when both spaces use the supremum norm, but it is unbounded.
\end{noteremark}

\begin{notetheorem}[Banach closed graph theorem]
Let $X$ and $Y$ be Banach spaces.  An everywhere-defined linear map
$T\colon X\to Y$ is bounded if and only if it is closed.
\end{notetheorem}

\begin{proof}
If $T$ is bounded and $x_n\to x$ while $Tx_n\to y$, then continuity gives
$Tx_n\to Tx$ and hence $y=Tx$.  Thus bounded operators are closed.

Conversely, suppose $T$ is closed and equip $X$ with the graph norm
\[
  \|x\|_T=\|x\|_X+\|Tx\|_Y.
\]
We first check the completeness assertion used in the source.  If $(x_n)$ is
Cauchy in $\|\cdot\|_T$, then $(x_n)$ is Cauchy in $X$ and $(Tx_n)$ is
Cauchy in $Y$.  Since both spaces are Banach, there are $x\in X$ and $y\in Y$
with
\[
  x_n\to x,
  \qquad
  Tx_n\to y.
\]
Closedness gives $y=Tx$, and therefore
\[
  \|x_n-x\|_T
  =\|x_n-x\|_X+\|Tx_n-Tx\|_Y\longrightarrow0.
\]
Hence $(X,\|\cdot\|_T)$ is Banach.

The identity map
\[
  I\colon(X,\|\cdot\|_T)\longrightarrow(X,\|\cdot\|_X)
\]
is a bounded bijection because $\|x\|_X\le\|x\|_T$.  By the bounded inverse
theorem, its inverse is bounded: there is $C>0$ such that
\[
  \|x\|_T\le C\|x\|_X
  \qquad(x\in X).
\]
In particular,
\[
  \|Tx\|_Y\le\|x\|_T\le C\|x\|_X,
\]
so $T$ is bounded.
\end{proof}

\begin{notetheorem}[Everywhere-defined symmetric maps]
Let $H$ be a real Hilbert space and let $T\colon H\to H$ be a map satisfying
\[
  \inner{Tx}{y}=\inner{x}{Ty}
  \qquad(x,y\in H).
\]
Then $T$ is linear and bounded.  The same conclusion holds over a complex Hilbert space when the inner
product is taken linear in the first variable and the symmetric identity
$\inner{Tx}{y}=\inner{x}{Ty}$ is imposed.
\end{notetheorem}

\begin{proof}
We first prove linearity without assuming continuity.  For
$x_1,x_2,y\in H$,
\[
\begin{aligned}
  \inner{T(x_1+x_2)-Tx_1-Tx_2}{y}
   &=\inner{x_1+x_2}{Ty}
     -\inner{x_1}{Ty}
     -\inner{x_2}{Ty}\\
   &=0.
\end{aligned}
\]
Since a vector orthogonal to every $y$ is zero,
$T(x_1+x_2)=Tx_1+Tx_2$.  If $\alpha$ is a scalar and the inner product is
linear in the first variable, then
\[
\begin{aligned}
  \inner{T(\alpha x)-\alpha Tx}{y}
   &=\inner{\alpha x}{Ty}-\alpha\inner{x}{Ty}=0
   \qquad(y\in H).
\end{aligned}
\]
Thus $T(\alpha x)=\alpha Tx$, and $T$ is linear.  With the opposite
inner-product convention the same computation is made in the other variable.

Now suppose $x_n\to x$ and $Tx_n\to z$.  For each $y\in H$,
\[
\begin{aligned}
  \inner{z}{y}
   &=\lim_{n\to\infty}\inner{Tx_n}{y}\\
   &=\lim_{n\to\infty}\inner{x_n}{Ty}\\
   &=\inner{x}{Ty}
    =\inner{Tx}{y}.
\end{aligned}
\]
Hence $z=Tx$, so $T$ is closed.  The closed graph theorem now applies to the
everywhere-defined linear operator $T$ and shows that $T$ is bounded.
\end{proof}

\section{Adjoints, compact maps, and closed ranges}

\begin{notedefinition}
For normed spaces $X,Y$ and $T\in\Bop{X}{Y}$, the adjoint
$T^*\colon Y^*\to X^*$ is defined by
\[
  (T^*y^*)(x)=y^*(Tx).
\]
A bounded linear map $T\colon X\to Y$ is \emph{compact} if it maps bounded
sets to relatively compact sets, equivalently, if every bounded sequence
$(x_n)$ has a subsequence for which $(Tx_{n_k})$ converges in $Y$.
\end{notedefinition}

\begin{notetheorem}[Schauder theorem]
A bounded linear operator $T\colon X\to Y$ is compact if and only if its
adjoint $T^*\colon Y^*\to X^*$ is compact.
\end{notetheorem}

\begin{noteremark}
A proof is given in Chapter~6, where compact operators are studied
systematically.
\end{noteremark}

\begin{notetheorem}[Closed range theorem]
Let $X,Y$ be Banach spaces and $T\in\Bop{X}{Y}$.  Then the following are
equivalent:
\begin{enumerate}
  \item $\Ran(T)$ is norm closed in $Y$;
  \item $\Ran(T^*)$ is norm closed in $X^*$;
  \item $\Ran(T^*)$ is weak-$*$ closed in $X^*$.
\end{enumerate}
Whenever these conditions hold,
\[
  \Ran(T)=\preann{\Ker(T^*)},
  \qquad
  \Ran(T^*)=\ann{\Ker(T)}.
\]
Without a closed-range hypothesis one always has
\[
  \overline{\Ran(T)}=\preann{\Ker(T^*)},
  \qquad
  \overline{\Ran(T^*)}^{\,w^*}=\ann{\Ker(T)}.
\]
\end{notetheorem}

\begin{notecorollary}
For $T\in\Bop{X}{Y}$ between Banach spaces, the following are equivalent:
\begin{enumerate}
  \item $T$ is surjective;
  \item there is $c>0$ such that $\|T^*y^*\|\ge c\|y^*\|$ for every
        $y^*\in Y^*$;
  \item $T^*$ is injective and has closed range.
\end{enumerate}
\end{notecorollary}

\begin{proof}
If $T$ is surjective, its induced map $X/\Ker T\to Y$ is a bounded bijection,
so the bounded inverse theorem yields the lower bound on $T^*$.  Conversely,
the lower bound makes $T^*$ injective with closed range.  The closed range
theorem then gives
$\overline{\Ran(T)}=\preann{\Ker(T^*)}=Y$, while $\Ran(T)$ is closed, so
$\Ran(T)=Y$.
\end{proof}
